\section{Slashing-Safe Fault Attribution}

\label{sec:slashing}

\subsection{Problem}

Proof-of-stake protocols rely on slashing to deter misbehavior: if a validator signs
conflicting blocks, their stake is destroyed. With threshold signatures, the entity
producing the signature is a \emph{group}, not an individual. How do we attribute
fault?

\subsection{CGGMP21 Identifiable Abort}

CGGMP21's identifiable abort property directly addresses this: if the signing protocol
produces an invalid output or fails, the honest parties can identify the corrupted
party. The proof of misbehavior is a verifiable transcript of the protocol messages
that demonstrates the corrupted party deviated from the protocol specification.

\subsection{Threshold Slashing Protocol}

\begin{enumerate}[leftmargin=2em]
  \item Each signing session produces a transcript $T = \{(i, m_i, \pi_i)\}$ where
        $m_i$ is party $i$'s message and $\pi_i$ is a zero-knowledge proof of
        correctness.
  \item If the final signature is invalid, the coordinator replays the transcript
        and identifies the first party whose message fails verification.
  \item A \emph{slashing proof} $\pi_{\text{slash}}$ is submitted on-chain:
        the party's message, the verification failure, and the transcript prefix
        proving all prior messages were valid.
  \item The on-chain slashing contract verifies $\pi_{\text{slash}}$ and destroys
        the identified party's stake.
\end{enumerate}

\begin{proposition}[Slashing soundness]
An honest threshold signer is never falsely slashed, assuming the discrete logarithm
assumption holds and the zero-knowledge proofs are sound.
\end{proposition}

%% ========================================================================
%%  8. APPLICATION TO LUX
%% ========================================================================
